\section{总结与延伸}

这堂课用一条连续谱组织了 retrieval-augmented language models：从 frozen BM25/vector search，把文档拼进 prompt；到 RePlug/reranker 让 retrieval 适配 frozen LM；再到 RAG、REALM、Atlas 与 RETRO 让 task signal 进入更多组件；最后把 retrieval 推广为 active tool use、multimodal memory 与 systems-over-models。

\subsection{统一比较框架}

任何 RAG 系统都可以用五元组描述：corpus/index $I$、retrieval policy $\pi$、retriever parameters $\theta_R$、fusion mechanism $F$、generator parameters $\theta_G$。比较论文或产品时，应报告这五项的训练状态、部署状态与成本，而不是只说“用了 RAG”。

\begin{importantbox}{本讲的十条可执行结论}
\begin{enumerate}
  \item 把 hallucination、generic error 与 evidence contradiction 分开评测；
  \item 同时报告 retrieval recall、evidence position、utilization 与 final correctness；
  \item exact entity 优先保留 sparse signal，语义召回再用 dense/late interaction；
  \item Frozen RAG 是强 baseline，但要记录 score/objective mismatch；
  \item 选择 RePlug、reranker 或 joint training 取决于 generator access level；
  \item 报告哪些参数更新、index 何时刷新、corpus 是否变化；
  \item 控制 $k_r$、$k_c$ 与 context position，避免盲目扩大 prompt；
  \item Active retrieval 要优化 information value 与 retrieval cost；
  \item Fine-tuning 改行为与表示，retrieval 改 per-request evidence，两者可组合；
  \item Source policy、provenance、permissions 与 evaluation 属于 RAG architecture。
\end{enumerate}
\end{importantbox}

\subsection{复现实验时应该记录什么}

至少记录 corpus snapshot、license/provenance、chunk size/overlap、embedding model、index type、retriever top-$k$、reranker、context packing、generator version、prompt、updated parameters、index refresh cadence、latency/cost 和 evaluation set。缺少这些信息，即使模型名称相同也难以复现。

\begin{warningbox}{不要把 2023 年课堂快照写成 2026 年产品事实}
本讲对 off-the-shelf retriever、vector database、hardware 与产业趋势的评价发生在 2023-12-05。讲义保留机制与历史判断，不据此断言 2026 年的当前产品能力、市场格局或法律结论。
\end{warningbox}

\subsection{拓展阅读}

建议按机制顺序阅读：DPR/ColBERT 理解 retrieval representation；RAG/FiD 理解 evidence marginalization 与 fusion；kNN-LM/RETRO 理解 non-parametric memory 如何进入 language modeling；REALM/Atlas 理解 joint training 的近似；RePlug/In-Context RALM 理解 frozen generator；FLARE/Self-RAG/Toolformer 理解 active retrieval 与 tool use；SILO/Lost-in-the-Middle 理解数据边界与 evidence utilization。

\subsection{最终小结}

Retrieval augmentation 的价值不只是“给模型更多文本”，而是把知识状态、更新、来源与计算从单一参数模型中拆出来。高质量 RAG 必须让 retriever、context、generator 与 source policy 对齐，并用分层 evaluation 证明：正确证据被找到、被放在可用位置、被生成器真正使用，最终输出也满足任务和 ground truth。

\end{document}
